% Target length ~150 words. Numbers fill via scripts/fill_results.py once runs land.
Catastrophic forgetting remains a central obstacle when adapting large language models through sequential fine-tuning. While parameter-efficient fine-tuning (PEFT) methods such as LoRA \citep{hu2022lora} mitigate forgetting at the parameter-count level, they treat every transformer component as equally susceptible. We ask a structural question: \textit{which components forget, and which resist forgetting?} We present the first per-component causally-traced \textbf{forgetting topology} for Qwen3-8B, an open 8B-scale model, mapping the forgetting susceptibility of every attention head and MLP block across a four-task sequence (NER $\rightarrow$ QA $\rightarrow$ Summarization $\rightarrow$ Code). We introduce the \textbf{Forgetting Vulnerability Score} (\fvs{}), a transferable metric computed as the weighted average of per-component causal-tracing deltas across task pairs, normalised to $[0,1]$, and \textbf{\tglora{}}, which restricts LoRA adapters to the $50\%$ most resistant components. \tglora{} matches standard LoRA on downstream tasks (within \pend{tg_lora_vs_lora_current}) while reducing prior-task forgetting by \pend{forget_reduction_pct} on average. The resulting topology yields a mechanistic explanation for \textit{why} fine-tuning hurts prior capabilities, and \textit{where}.
